\documentclass[11pt]{article}
\usepackage{quanttheme}

\newcolumntype{Y}{>{\raggedright\arraybackslash}X}
\newcommand{\tk}[1]{\texttt{#1}}

\begin{document}
\thispagestyle{empty}

\titlebanner
  {QuantAnomalies}
  {Quantitative Anomalies as Market Opportunities}
  {Progress Presentation \quad\textbullet\quad Statistical \& AI Dashboard for Equity and Forex Markets}

\vspace{1.2em}

\section{Title of the Study}
\textbf{Quantitative Anomalies as Market Opportunities — A Statistical \& AI
Dashboard for Equity and Forex Markets.}

\begin{opp}
In plain terms: textbook finance says markets are ``efficient'' — prices already
reflect everything, risk never changes, and nothing is predictable. Reality is
messier, and those cracks in the theory are recurring patterns called
\textbf{anomalies}. Our project hunts for those anomalies and reframes them as
\textbf{opportunities}: places where the odds may tilt in a careful observer's favour.
\end{opp}

\section{Objectives of the Study}
Each objective targets one well-known market anomaly and turns it into a usable
opportunity signal.

\begin{enumerate}
\item \textbf{Macro-driver opportunity.} Measure \emph{which} macroeconomic forces
(interest rates, inflation, the fear index, oil, gold, the dollar) move an asset's
returns — and crucially with \emph{what time delay}. If a driver acts weeks later,
that lag is a window to anticipate moves.

\item \textbf{Risk-regime opportunity.} Model how risk itself rises and falls over
time (volatility clustering) and detect that crashes happen far more often than a
bell curve predicts (fat tails). Spotting a ``stormy'' regime early is a chance to
cut risk — or to recognise unusually cheap calm.

\item \textbf{Mean-reversion opportunity.} Find pairs of currencies whose prices are
tethered over the long run (cointegration), then flag when they drift unusually far
apart — history says the gap tends to snap back. This is the basis of
market-neutral ``statistical arbitrage''.

\item \textbf{Opportunity synthesis (the decision engine).} Combine all of the above
into one automated scanner that \emph{fuses} what the four pillars detect into a
single directional verdict — with agreement, conviction, and disagreement all made
explicit — and uses a \textbf{local AI analyst} to explain it in plain English.

\item \textbf{Accessibility \& robustness.} Make it work for \emph{any} ticker the
user types, on a reliable data pipeline, presented through a clear, interactive
dashboard.
\end{enumerate}

\section{Motivation for Selecting the Topic}
\begin{itemize}
\item \textbf{The anomalies are real and well-studied, but rarely made tangible.}
Lagged macro effects, volatility clustering, fat tails, and cointegration all appear
in finance literature; we wanted to make them \emph{visible and actionable} rather
than abstract equations.
\item \textbf{Reframing risk as opportunity.} The same statistics used to warn about
risk can also point to opportunity — we wanted a tool that does both.
\item \textbf{Bridging statistics and modern AI.} We combine classic econometrics
(regression, GARCH, cointegration) with a \textbf{local large language model} that
turns numbers into a readable analyst note — private, free to run, and on our own GPU.
\item \textbf{Learning value.} The project spans real statistical modelling, live
financial-data engineering, full-stack web development, and on-device AI.
\end{itemize}

\section{Data \& APIs Used}
\textit{In place of a questionnaire and a survey sample: our ``instrument'' is a set
of financial-data APIs, and our ``sample'' is the live market history they return.}

\subsection{A. External data-source APIs (collecting raw market data)}
\begin{itemize}
\item \textbf{Yahoo Finance} (via the \tk{yfinance} library) — daily price history
(2015--2025) for any equity, index, currency pair, future, or crypto. Examples:
\tk{\textasciicircum{}GSPC} (S\&P 500), \tk{\textasciicircum{}VIX} (volatility index),
\tk{CL=F} (oil), \tk{GC=F} (gold), \tk{\textasciicircum{}TNX} (10-year Treasury
yield), \tk{DX-Y.NYB} (US dollar index), and 45 forex pairs.
\item \textbf{FRED — U.S. Federal Reserve Economic Data} — macroeconomic series:
CPI (inflation), Fed Funds Rate, and Unemployment. A free mirror, \textbf{DBnomics},
is a built-in backup so a single outage never stops the study.
\end{itemize}

\subsection{B. Our own analysis API (FastAPI backend)}
Turns raw data into results:
\begin{itemize}
\item \tk{/api/macro-regression/*} — regression, Granger-causality, correlation, time series
\item \tk{/api/garch/*} — volatility model, clustering tests, return distribution
\item \tk{/api/pairs/*} — cointegration tests, best-pair spread \& signals, correlation
\item \tk{/api/options/black-scholes} — option price, Greeks, implied volatility
\item \tk{/api/engine/\{feed,asset,narrate,status\}} — the fused opportunity scan, the
streaming AI narration, and the engine's health
\item \tk{/api/llm/info} — status of the local AI analyst
\end{itemize}

\subsection{C. Local AI API}
An on-device \tk{llama.cpp} server (OpenAI-compatible) running the \textbf{Qwen3-4B}
model on our \textbf{NVIDIA RTX 4050 GPU}, used to explain detected opportunities in
plain language. No data leaves the machine.

\section{Pilot Study Conducted to Date}
We have a \textbf{working end-to-end prototype}, validated on real data. \emph{All
figures below were recomputed on 26 July 2026.}

\begin{keybox}
\textbf{Data pipeline — and one real finding.} We had been treating a data corruption
as a bad cache. The actual root cause is that \textbf{\tk{yfinance} is not
thread-safe}: concurrent downloads silently return one ticker's data under another's
name. A deliberately threaded scan reproduced it on demand — three separate ticker
pairs came back byte-identical. Downloads are now serialized. Separately, the data
end-date had been frozen, so the dashboard was reasoning over seven-month-old prices;
now a rolling window with a staleness check.
\end{keybox}

\textbf{Three analysis pillars — working, with sample findings on the S\&P 500:}
\begin{itemize}
\item \emph{Macro:} standardized macro factors explain about \textbf{69\%} of monthly
return variation (R² = 0.687, adj.\ 0.585); Granger causality runs 32 tests,
\textbf{6 significant}. Coefficients are standardized betas, so ``which driver matters
most'' is a fair comparison.
\item \emph{Volatility:} GARCH persistence ≈ \textbf{0.994} over 2{,}905 days (shocks
fade very slowly); excess kurtosis ≈ \textbf{15.8}, skew \textbf{−0.65}, Jarque-Bera
p ≈ 0 — risk is \emph{not} constant, and the normal curve badly under-states crashes.
\item \emph{Pairs:} \textbf{USDCHF/USDJPY} tested cointegrated (p = \textbf{0.0075})
with a mean-reversion half-life of \textbf{69 days}, on log prices.
\end{itemize}

\textbf{Decision engine — the headline advance.} Eight detectors across all four
pillars, normalized onto one bull/bear axis and then \textbf{fused}: weight =
severity × statistical reliability, producing a direction (\tk{tilt}) and a
\tk{conviction} that \emph{falls} when signals disagree, with risk on a separate axis
so a volatility spike cannot fake a directional call. Signals that oppose the verdict
are shown, not averaged away. Warm response: \textbf{11\,ms} for the cross-asset feed,
\textbf{4.5\,ms} per asset.

\textbf{The AI layer, kept honest.} The model is handed the computed numbers and
explains them; it never computes. Its own stance renders \emph{beside} the computed
one, so a disagreement is visible. Every figure in its note is checked against the
supplied evidence and anything unmatched is flagged in the UI. ≈4 seconds on the GPU
(≈28 tokens/sec), streaming, fully local.

\textbf{Dashboard (working).} Any ticker, any forex pairs, time ranges on every chart,
six tabs (Overview, Opportunities, Macro, GARCH, Pair Trading, Options), light/dark
themes
with all colour pairings verified at ≥4.5:1 contrast.

\begin{opp}
\textbf{In short:} the pilot confirms the data, the statistics, the fusion layer, the
AI explanation, and the interface all work together on live markets. The honest
remaining gap is \textbf{inferential rigor} — multiple-testing correction on the 990
cointegration tests, HAC standard errors on the autocorrelated monthly regression, and
an out-of-sample check on the signals. That is the next block of work.
\end{opp}

\section{Modules Built So Far}

\subsection*{Backend (Python / FastAPI)}
\begin{itemize}
\item \tk{data\_loader.py} — data acquisition (Yahoo Finance, FRED, DBnomics), caching, dataset builders
\item \tk{analysis/macro\_regression.py} — Pillar 1: OLS lag regression, Granger causality, correlation
\item \tk{analysis/garch.py} — Pillar 2: GARCH(1,1), volatility clustering, return distribution
\item \tk{analysis/pairs.py} — Pillar 3: cointegration, spread / z-score, trading signals
\item \tk{analysis/recommender.py} — the eight detectors over the pillar outputs
\item \tk{analysis/black\_scholes.py} — Pillar 4: option pricing, Greeks, implied-volatility solver
\item \tk{engine.py} — the decision engine: polarity normalisation, reliability
weighting, fusion, the tiered/cached scan, and streaming narration
\item \tk{llm\_client.py} — provider-agnostic local-LLM client (GPU server + CPU fallback)
\item \tk{main.py} — FastAPI endpoints tying it all together
\item \tk{test\_engine.py} — 25 tests over polarity, fusion invariants, ranking,
reliability, stance parsing, and the number guardrail
\end{itemize}

\subsection*{Frontend (React + Vite)}
\begin{itemize}
\item Six tab panels — Overview, Opportunities, Macro Regression, GARCH Volatility, Pair Trading, Options
\item Shared components — SVG icon set, info-tooltips, time-range filter, ticker
search, forex selector, status states, tilt gauge, skeletons
\item Support — design tokens, light/dark theme hook, shared ticker context,
SSE narration hook, number formatting, data-fetch hook with abort + session cache
\end{itemize}

\section{Future Plans — Expanding the Opportunity Engine}
Each future module is another \emph{lens} for spotting opportunity. The engine's
principle stays the same: \textbf{statistics detect, the AI explains}.

\subsection*{1. Options \& Black–Scholes \quad{\normalfont\small(delivered; volatility forecast still open)}}
Price European call / put options and their Greeks, and back out the market's
\textbf{implied volatility}. Now shipped end to end: the \textbf{Options} tab prices
any contract on a chosen strike and expiry, and \tk{detect\_vol\_mispricing} feeds the
engine as its eighth detector. \textbf{Opportunity:} when the market pays far more (or
less) for risk than the asset has actually delivered, options are ``rich'' or
``cheap''. \emph{Remaining:} the comparison currently uses 1-year realised volatility;
driving it off a GARCH forecast over the option's own tenor is the open piece.

\subsection*{2. Valuation — ``is the stock bloated / overpriced?''}
Pull fundamental ratios (P/E, P/B, P/S, PEG, EV/EBITDA) and compare them with the
stock's own history and its peers, alongside a technical ``stretch'' check (how far
price sits above its 200-day average, and RSI). \textbf{Opportunity:} flag names
``priced for perfection'' (avoid / potential short) versus those ``on sale''
(potential value buys).

\subsection*{3. More opportunity lenses}
\begin{itemize}
\item \textbf{Momentum \& relative strength} — rank a universe by risk-adjusted
12–1 month momentum to spot rotation leaders and laggards.
\item \textbf{Oversold bounce} — distance from the 52-week high plus RSI to find
mean-reversion candidates.
\item \textbf{Volatility squeeze} — detect compression that often precedes a breakout.
\item \textbf{Correlation-regime shift} — assets that usually move together
decoupling (a diversification or pairs opportunity).
\item \textbf{Seasonality / calendar effects} — turn-of-month and ``sell in May''
style documented anomalies.
\item \textbf{Alpha vs beta} — separate market-driven moves from stock-specific mispricing.
\item \textbf{Unusual volume} — accumulation / distribution versus the average.
\item \textbf{Risk budgeting} — turn the detected regime into a suggested position size.
\end{itemize}

\section{Decision Engine: Rules-based vs LLM-based}
We combine them: \emph{rules decide, the LLM explains}. Rules alone is the default and
works with no model running. Where the two meet, the design keeps the seam visible —
the model's stance renders beside the computed stance so it may disagree openly, and
every figure it writes is checked against the evidence it was given.

\begin{table}[h]
\renewcommand{\arraystretch}{1.3}
\begin{tabularx}{\linewidth}{>{\bfseries}l Y Y}
\toprule
& \textbf{\color{qnavy}Rules-based} & \textbf{\color{qnavy}LLM-based} \\
\midrule
Determinism & Fully deterministic, reproducible & Varies with settings \\
Explainability & Transparent thresholds & Natural-language reasoning \\
Numbers & Exact by construction & Must be guarded against errors \\
Flexibility & Rigid; new rules need code & Synthesises novel combinations \\
Speed \& cost & Instant, free & ≈4 s on GPU, compute cost \\
Works offline & Yes & Needs the local model \\
Best at & Detection \& signals & Explanation \& synthesis \\
\bottomrule
\end{tabularx}
\end{table}

\begin{opp}
Our design uses each for what it is good at — rules produce the trustworthy numbers,
the LLM produces the readable story. A future \textbf{LLM decision-maker} mode would
let the model also weigh and rank signals (not just explain), kept honest by
rule-based guardrails and number validation.
\end{opp}

\section*{Appendix — Each anomaly in one line (layman)}
\begin{itemize}
\item \textbf{Lagged macro effect:} \textit{News ripples through markets over weeks,
not instantly — know the lever and the delay, and you can position ahead.}
\item \textbf{Volatility clustering:} \textit{Calm days cluster, and so do wild days —
storms arrive in groups, giving you warning to manage risk.}
\item \textbf{Fat tails:} \textit{Extreme crashes happen much more often than the
textbook bell curve says — plan for them.}
\item \textbf{Cointegration / pairs:} \textit{Two currencies on one leash: when they
drift far apart, they usually snap back — bet on the gap closing, up market or down.}
\end{itemize}

\end{document}
